\lecture{19}{10. November 2025}{Superposition}

\begin{problemwithimage}{f16_38}{16-38}
  Determine the vertical deflection and slope at the end $A$ of the bracket.
  Assume that the bracket is fixed supported at its base, and neglect the axial deformation of segment $AB$.
  $EI$ is constant.
\end{problemwithimage}

\begin{derivation}
  We split the problem into two cases, the distributed loading on $AB$ and the point loading $P$ on $BC$ which will then be added at the end using the principle of superposition.

  \paragraph{Distributed loading case}
  For the cantilevered beam $AB$ we have the maximum deflection and slope from Appendix D of the book as:
  \begin{equation*}
    v_{\mathrm{max}}^{w} = -\frac{wL^{4}}{8EI} \qquad \theta^{w}_{\mathrm{max}} = - \frac{wL^3}{6EI}
  .\end{equation*}
  The moment at $B$ from the distributed loading is:
  \begin{equation*}
    M_0^{w} = \frac{w}{2} \cdot L^2
  .\end{equation*}
  
  This moment will deflect the $BC$ cantilever beam too. From Appendix D we see that the total slope in $BC$ caused by the moment $M_0$ applied at the free end is:
  \begin{equation*}
    \theta_{B}^{w} = \frac{M_0 h}{EI} = \frac{w L^2 h}{2EI}
  .\end{equation*}
  This corresponds to a deflection at $A$ of
  \begin{equation*}
    v_A^{w} = L \theta_B^{w} = L \frac{wL^2 h}{2EI} = \frac{wL^3h}{2EI}
  .\end{equation*}
  That aforementioned rotation also adds to the slope at $A$:
  \begin{equation*}
    \theta_A^{w} = \frac{wL^2 h}{2EI}
  .\end{equation*}
  We add this to the standard tabulated formulas to get:
  \begin{align*}
    v_A^{(1)} &= v_{\mathrm{max}}^{w} + v_A^{w} = - \frac{wL^3h}{2EI} - \frac{2L^{4}}{8EI} \\
    \theta_A^{(1)} &= \theta_{\mathrm{max}}^{w} + \theta_A^{w} = - \frac{wL^2h}{2EI} - \frac{wL^3}{6EI}
  .\end{align*}
  
  \paragraph{Point loading case}
  We treat $BC$ as a cantilever with a load $P = \qty{400}{\N}$ applied to the tip.
  This gives the standard formula for the deflection:
  \begin{equation*}
    \theta_B^{P} = - \frac{PL^2}{2EI}
  .\end{equation*}
  Which swings $A$ down by
  \begin{equation*}
    v_A^{p} = \theta_B^{p} \cdot h = - \frac{PL^2h}{2EI}
  .\end{equation*}
  And the slope at $A$ from this case is just that rigid rotation:
  \begin{equation*}
    \theta_A^{p} = - \frac{PL^2}{2EI}
  .\end{equation*}
  
  \paragraph{Superposition}
  Now we just add those results:
  \begin{align*}
    v_A &= v_A^{(1)} + v_A^{p} = - \frac{wL^3h}{2EI} - \frac{2L^{4}}{8EI} - \frac{PL^2h}{2EI} \\
    \theta_{A} &= \theta_A^{(1)} + \theta_A^{p} = - \frac{wL^2h}{2EI} - \frac{wL^3}{6EI} - \frac{PL^2}{2EI}
  .\end{align*}
  Or numerically:
  \begin{align*}
    v_A &= - \qty{312,5e6}{\mm} \cdot \frac{1}{EI} \\
    \theta_A &= - \qty{3,2917e6}{rad} \cdot \frac{1}{EI}
  .\end{align*}
\end{derivation}

\finalresult{
\begin{aligned}
  v_A &= - \qty{312,5e6}{\mm} \cdot \frac{1}{EI} \\
  \theta_A &= - \qty{3,2917e6}{rad} \cdot \frac{1}{EI}
.\end{aligned}
}

\begin{problemwithimage}{f16_43}{16-43}
  The simply supported beam carries a uniform load of \qty{30}{\kN\per\m}.
  Code restrictions, due to a plaster ceiling, require the maximum deflection not to exceed $1 / 360$ of the span length.
  Select the lightest-weight A--36 steel wide-flange beam from Appendix C that will satisfy this requirement and safely support the load.
  The allowable bending stress is $\sigma_{\mathrm{allow}} = \qty{168}{\MPa} $ and the allowable shear stress is $\tau_{\mathrm{allow}} = \qty{100}{\MPa}$.
  Assume $A$ is a pin and $B$ a roller support.
\end{problemwithimage}

\begin{derivation}
  We start by converting the distributed load into an equivalent point loading as
  \begin{equation*}
    W = wL =  \qty{30}{\kN\per\m} \cdot \qty{4.8}{\m}  = \qty{144}{\kN}
  .\end{equation*}
  A force equilibrium gives:
  \begin{align*}
    \sum F_y &= 0 \\
    A_y + B_y - \qty{40}{\kN} - \qty{144}{\kN} - \qty{40}{\kN} &= 0 \\
    A_y + B_y &= \qty{224}{\kN} 
  .\end{align*}
  Due to symmetry $A_y = B_y = \qty{112}{\kN}$ so
  \begin{equation*}
    V_{\mathrm{max}} = \qty{112}{\kN}
  .\end{equation*}
  
  From symmetry, the maximum moment is in the middle (where shear crosses zero).
  The moment at the middle is:
  \begin{equation*}
    M_{\mathrm{max}} = R_A x - \frac{wx^2}{2}  - P \left( x - \qty{1,2}{\m}  \right) = \qty{112}{\kN} \cdot \qty{2.4}{\m} - \frac{\qty{30}{\kN\per\m} \cdot \left( \qty{2.4}{\m}  \right)^2}{2} - \qty{40}{\kN} \left( \qty{2.4}{\m} - \qty{1.2}{\m}  \right) = \qty{134,4}{\kN.\m}
  .\end{equation*}
  
  The allowable bending stress is $\sigma_{\mathrm{allow}} = \qty{168}{\MPa}$. 
  This means the elastic section modulus $S$ has to be at least:
  \begin{equation*}
    S_{x, \mathrm{min}} = \frac{M_{\mathrm{max}}}{\sigma_{\mathrm{allow}}} = \frac{\qty{134.4}{\kN.\m}}{\qty{168}{\MPa}} = \qty{8e5}{\mm^3} = \qty{48.82}{in^3}  
  .\end{equation*}
  
  The deflection requirement means the maximum allowed deflection is:
  \begin{equation*}
    \delta_{\mathrm{max}} = \frac{1}{360} \cdot \qty{4.8}{\m}  = \qty{13.33}{mm}
  .\end{equation*}
  
  We now use superposition.
  The deflection from the distributed load is:
  \begin{equation*}
    \delta_w = \frac{5wL^{4}}{384EI}
  .\end{equation*}
  And the deflection from the point loads are
  \begin{equation*}
    \delta_{2P} = \frac{Pa \left( \frac{3}{4}L^2 - a^2 \right)}{6EI}
  .\end{equation*}
  So
  \begin{align*}
    \delta &= \delta_w + \delta_{2P} \\
           &= \frac{5 w L^{4}}{384 EI} + \frac{Pa \left( \frac{3}{4} L^2 - a^2 \right)}{6EI} = \frac{5 \cdot \qty{30}{\kN\per\m} \cdot \left( \qty{4.8}{\m}  \right)^{4}}{384 \cdot \qty{200}{\GPa}} \cdot \frac{1}{I} + \frac{\qty{40}{\kN} \cdot \qty{1.2}{\m} \left( \frac{3}{4} \cdot \left( \qty{4.8}{\m}  \right)^2 - \left( \qty{1.2}{\m}  \right)^2 \right)}{6\cdot \qty{200}{\GPa} } \cdot \frac{1}{I} \\
           &= \qty{1.0368e-6}{\m^5} \cdot \frac{1}{I} + \qty{6.336e-7}{\m^5} \cdot \frac{1}{I} \\
           &= \frac{\qty{6,336e-7}{\m^{5}}}{I}
  .\end{align*}
  Which means:
  \begin{equation*}
    I_{\mathrm{min}} = \frac{\qty{6,336e-7}{\m^{5}}}{\qty{13,33}{\mm} } = \qty{301}{in ^{4}} 
  .\end{equation*}
  
  The lightest steel satisfying this in Appendix C is W14x38.
\end{derivation}

\newpage
